\chapter{Registers}
\label{ch:registers}

\section{Register Map}

The registers occupy the first words of the device segment. Every defined register is 16 bits wide and is read and
written a word at a time.

\begin{table}[H]
	\centering
	\begin{tabular}{lllp{5.2cm}}
		\toprule
		\textbf{Offset}                  & \textbf{Name}     & \textbf{Access} & \textbf{Description}                                     \\
		\midrule
		\texttt{0x0000}                  & \texttt{CONFIG}   & R/W             & Base configuration: mode, layer disables, brightness     \\
		\texttt{0x0001}                  & \texttt{TILESIZE} & R/W             & Tile and tilemap sizes                                   \\
		\texttt{0x0002}                  & \texttt{B1MAP}    & R/W             & Background 1 tilemap address                             \\
		\texttt{0x0003}                  & \texttt{B2MAP}    & R/W             & Background 2 tilemap address                             \\
		\texttt{0x0004}                  & \texttt{B3MAP}    & R/W             & Background 3 tilemap address                             \\
		\texttt{0x0005}                  & ---               & Reserved        &                                                          \\
		\texttt{0x0006}                  & \texttt{B1TILES}  & R/W             & Background 1 tile data address                           \\
		\texttt{0x0007}                  & \texttt{B2TILES}  & R/W             & Background 2 tile data address                           \\
		\texttt{0x0008}                  & \texttt{B3TILES}  & R/W             & Background 3 tile data address                           \\
		\texttt{0x0009}                  & ---               & Reserved        &                                                          \\
		\texttt{0x000A}                  & \texttt{B1SCRX}   & R/W             & Background 1 horizontal scroll                           \\
		\texttt{0x000B}                  & \texttt{B2SCRX}   & R/W             & Background 2 horizontal scroll                           \\
		\texttt{0x000C}                  & \texttt{B3SCRX}   & R/W             & Background 3 horizontal scroll                           \\
		\texttt{0x000D}                  & ---               & Reserved        &                                                          \\
		\texttt{0x000E}                  & \texttt{B1SCRY}   & R/W             & Background 1 vertical scroll                             \\
		\texttt{0x000F}                  & \texttt{B2SCRY}   & R/W             & Background 2 vertical scroll                             \\
		\texttt{0x0010}                  & \texttt{B3SCRY}   & R/W             & Background 3 vertical scroll                             \\
		\texttt{0x0011}                  & \texttt{INTCTL}   & R/W             & Interrupt enables                                        \\
		\texttt{0x0012}                  & \texttt{B1PAL}    & R/W             & Background 1 palette configuration                       \\
		\texttt{0x0013}                  & \texttt{B2PAL}    & R/W             & Background 2 palette configuration                       \\
		\texttt{0x0014}                  & \texttt{B3PAL}    & R/W             & Background 3 palette configuration                       \\
		\texttt{0x0015}                  & \texttt{LINE}     & R/W             & Line-compare interrupt line; a write acknowledges        \\
		\texttt{0x0016}                  & \texttt{STILES}   & R/W             & Sprite tile data address                                 \\
		\texttt{0x0017}                  & \texttt{STATUS}   & R               & Version, output mode, sprite overflow, vertical blanking \\
		\texttt{0x0018}                  & \texttt{VLINE}    & R               & Current line                                             \\
		\texttt{0x0019}--\texttt{0x00FF} & ---               & Open bus        & Reads zero, writes ignored                               \\
		\bottomrule
	\end{tabular}
	\caption{Register Map}
	\label{tab:register-map}
\end{table}

The register names are those of this manual. Software must not write the reserved registers. The registers are laid out
in groups of three backgrounds followed by one reserved slot, so a register for background $n$ is at the group's base
plus $n - 1$.

\section{Base Configuration Register}
\label{sec:reg-base-config}

\begin{figure}[H]
	\centering
	\begin{bytefield}[bitwidth=2em]{16}
		\bitheader{0-15} \\
		\bitbox{1}{GD} &
		\bitbox{1}{B1D} &
		\bitbox{1}{B2D} &
		\bitbox{1}{B3D} &
		\bitbox{1}{} &
		\bitbox{1}{S1D} &
		\bitbox{1}{S2D} &
		\bitbox{1}{S3D} &
		\bitbox{1}{\footnotesize SHD} &
		\bitbox{4}{BRT} &
		\bitbox{3}{MODE}
	\end{bytefield}
	\caption{Base Configuration Register (\texttt{CONFIG}, Offset \texttt{0x0000})}
	\label{fig:reg-config}
\end{figure}

\begin{table}[H]
	\centering
	\begin{tabular}{llp{7.5cm}}
		\toprule
		\textbf{Bits} & \textbf{Field}                           & \textbf{Meaning}                                                                                                                                                                             \\
		\midrule
		15            & \texttt{GD}                              & Graphics disable (forced blank): when 1 the picture is black and the device makes no reads, so video memory and the sprite table can be accessed at any time. 1 at reset                     \\
		14, 13, 12    & \texttt{B1D}, \texttt{B2D}, \texttt{B3D} & Disable background 1, 2, 3 when 1 (Section~\ref{sec:m2-layers})                                                                                                                              \\
		11            & ---                                      & Reserved                                                                                                                                                                                     \\
		10, 9, 8      & \texttt{S1D}, \texttt{S2D}, \texttt{S3D} & Disable sprite layer 1, 2, 3 when 1 (Chapter~\ref{ch:sprites})                                                                                                                               \\
		7             & \texttt{SHD}                             & Sprite shadow enable: when 1, value 15 of sprite palette 7 is a shadow (Section~\ref{sec:sprite-shadow}); 0 at reset                                                                         \\
		6--3          & \texttt{BRT}                             & Screen brightness: each colour channel of the final picture is multiplied by $(\texttt{BRT} + 1)/16$, so 15 is full brightness and 0 nearly black. \texttt{GD} blanks the picture completely \\
		2--0          & \texttt{MODE}                            & Video mode: 0 for Mode 1, 1 for Mode 2, 2 for Mode 3, others reserved                                                                                                                        \\
		\bottomrule
	\end{tabular}
	\caption{Base Configuration Register Fields}
	\label{tab:reg-config-fields}
\end{table}

The \texttt{MODE} field is sampled at the first cycle of each line, so a change takes effect at the start of the next
line. Software should change mode during vertical blanking.

\section{Tile Size Register}
\label{sec:reg-tilesize}

\begin{figure}[H]
	\centering
	\begin{bytefield}[bitwidth=2em]{16}
		\bitheader{0-15} \\
		\bitbox{1}{B1S} &
		\bitbox{1}{B2S} &
		\bitbox{1}{B3S} &
		\bitbox{1}{} &
		\bitbox{1}{S1S} &
		\bitbox{1}{S2S} &
		\bitbox{1}{S3S} &
		\bitbox{1}{} &
		\bitbox{1}{B1X} &
		\bitbox{1}{B2X} &
		\bitbox{1}{B3X} &
		\bitbox{1}{} &
		\bitbox{1}{B1Y} &
		\bitbox{1}{B2Y} &
		\bitbox{1}{B3Y} &
		\bitbox{1}{}
	\end{bytefield}
	\caption{Tile Size Register (\texttt{TILESIZE}, Offset \texttt{0x0001})}
	\label{fig:reg-tilesize}
\end{figure}

\begin{table}[H]
	\centering
	\begin{tabular}{llp{7.5cm}}
		\toprule
		\textbf{Bits} & \textbf{Field}                           & \textbf{Meaning}                                                                          \\
		\midrule
		15, 14, 13    & \texttt{B1S}, \texttt{B2S}, \texttt{B3S} & Orthogonal background tile size: 0 for 8 by 8, 1 for 16 by 16 (Section~\ref{sec:flat-16}) \\
		11, 10, 9     & \texttt{S1S}, \texttt{S2S}, \texttt{S3S} & Sprite tile size, same encoding                                                           \\
		7, 6, 5       & \texttt{B1X}, \texttt{B2X}, \texttt{B3X} & Tilemap width: 0 for one screen of 32 entries, 1 for two (Section~\ref{sec:tilemap-size}) \\
		3, 2, 1       & \texttt{B1Y}, \texttt{B2Y}, \texttt{B3Y} & Tilemap height: 0 for one screen of 32 entries, 1 for two                                 \\
		12, 8, 4, 0   & ---                                      & Reserved                                                                                  \\
		\bottomrule
	\end{tabular}
	\caption{Tile Size Register Fields}
	\label{tab:reg-tilesize-fields}
\end{table}

An \emph{orthogonal} background is an ordinary background of square tiles in rows and columns, as opposed to an
isometric one. The word is used for all of them: the three backgrounds of Mode 1, the two of Mode 3 and the status
layer of Mode 2. Orthogonal backgrounds scroll, except for the status layer of Mode 2. The tile size bits
\texttt{B1S}--\texttt{B3S} apply to orthogonal backgrounds, that is, to all three backgrounds in Mode 1 and to
backgrounds 1 and 2 in Mode 3, but not to the status layer, which always uses 8 by 8 tiles. The tilemap size bits apply
to orthogonal and isometric backgrounds. Isometric backgrounds always use 16 by 16 texel tiles and ignore the tile size
bits (Section~\ref{sec:m2-restrictions}).

\section{Tilemap and Tile Data Address Registers}

\texttt{B1MAP}--\texttt{B3MAP} hold the word address of the first entry of the background's tilemap. \texttt{B1TILES}--\texttt{B3TILES}
hold the word address of the first word of tile number 0 of its tile data. Both are full offsets within the segment, so
a value for VRAM is at least \texttt{0x8000} (Chapter~\ref{ch:memory-map}). All 16 bits are significant. The renderer
adds the offset of an entry or tile to the base with 16-bit wraparound.

\texttt{STILES} holds the sprite tile data address (Chapter~\ref{ch:sprites}).

\section{Scroll Registers}

\begin{figure}[H]
	\centering
	\begin{bytefield}[bitwidth=2em]{16}
		\bitheader{0-15} \\
		\bitbox{6}{} &
		\bitbox{10}{SCROLL}
	\end{bytefield}
	\caption{Scroll Registers (\texttt{B1SCRX}--\texttt{B3SCRY})}
	\label{fig:reg-scroll}
\end{figure}

Bits 9--0 hold the scroll amount, 0 to 1023, in pixels. Bits 15--10 are reserved. The amount moves the background's
view of its map: increasing \texttt{SCRX} moves the picture to the left, and increasing \texttt{SCRY} moves it up.

Every background of every mode scrolls. The amounts are sampled once per line, as the chapter of each mode describes.

\section{Palette Configuration Registers}
\label{sec:reg-palette}

\begin{figure}[H]
	\centering
	\begin{bytefield}[bitwidth=2em]{16}
		\bitheader{0-15} \\
		\bitbox{8}{OFFSET} &
		\bitbox{6}{} &
		\bitbox{2}{SIZE}
	\end{bytefield}
	\caption{Palette Configuration Registers (\texttt{B1PAL}--\texttt{B3PAL})}
	\label{fig:reg-pal}
\end{figure}

\begin{table}[H]
	\centering
	\begin{tabular}{llp{7.5cm}}
		\toprule
		\textbf{Bits} & \textbf{Field}  & \textbf{Meaning}                                                       \\
		\midrule
		15--8         & \texttt{OFFSET} & Palette entry added to every colour address of the background          \\
		7--3          & ---             & Reserved                                                               \\
		2             & \texttt{T2}     & 2-bit tiles (Sections~\ref{sec:m2-two-bit} and~\ref{sec:flat-two-bit}) \\
		1--0          & \texttt{SIZE}   & Colours per palette: 0 for 4, 1 for 8, 2 for 16, 3 for 32              \\
		\bottomrule
	\end{tabular}
	\caption{Palette Configuration Register Fields}
	\label{tab:reg-pal-fields}
\end{table}

Section~\ref{sec:palette-select} explains how \texttt{OFFSET} and \texttt{SIZE} combine with a tilemap entry's palette
selection to choose a palette entry. 4-bit tile data uses sizes of 16 colours or fewer.

\section{Status Register}

\begin{figure}[H]
	\centering
	\begin{bytefield}[bitwidth=2em]{16}
		\bitheader{0-15} \\
		\bitbox{8}{} &
		\bitbox{1}{VBL} &
		\bitbox{1}{LIP} &
		\bitbox{1}{SOV} &
		\bitbox{1}{OUT} &
		\bitbox{4}{VER}
	\end{bytefield}
	\caption{Status Register (\texttt{STATUS}, Offset \texttt{0x0017})}
	\label{fig:reg-status}
\end{figure}

\texttt{VER}, bits 3--0, is the version of the device, which is 1. \texttt{OUT}, bit 4, is the output mode: 0 for NTSC and
1 for PAL. Both are fixed when the device is made. \texttt{SOV}, bit 5, is the sprite overflow flag: it is set when any
line of the frame had more than 32 sprites (Section~\ref{sec:sprite-budget}), and cleared at the start of line 0 of the
next frame, so a program reads it in vertical blanking. \texttt{LIP}, bit 6, is set while a line-compare interrupt is
pending (Section~\ref{sec:line-interrupt}). \texttt{VBL}, bit 7, is set while the beam is in vertical blanking, from the first cycle
of line 224 to the last cycle of the last line of the frame, and is clear in lines 0 to 223; it is the flag a program
polls to know that video memory is open (Section~\ref{sec:memory-vram-access}). Bits 15--8 are reserved. The register is read-only.

\section{Vertical Line Register}
\label{sec:reg-vline}

\texttt{VLINE}, offset \texttt{0x0018}, holds in bits 8--0 the number of the line the beam is on: 0 to 261 on NTSC and 0 to 311 on
PAL. It changes at the first cycle of each line, the cycle at which a line-compare interrupt for that line is raised, and
bits 15--9 read as zero. The register is read-only. A program reads it to find where in the frame it is, and
reads it twice, or compares it with \texttt{LINE}, when a line may have passed between two accesses.

\section{Line Compare Register}
\label{sec:reg-line}

\texttt{LINE} holds, in bits 8--0, the number of the line at which the line-compare interrupt fires when \texttt{LIE} is set
in the Interrupt Control register (Section~\ref{sec:line-interrupt}). Bits 15--9 are reserved. A write to \texttt{LINE}, of
any value, also acknowledges a pending line-compare interrupt.

\section{Interrupt Control Register}
\label{sec:reg-intctl}

\begin{figure}[H]
	\centering
	\begin{bytefield}[bitwidth=2em]{16}
		\bitheader{0-15} \\
		\bitbox{14}{} &
		\bitbox{1}{LIE} &
		\bitbox{1}{VIE}
	\end{bytefield}
	\caption{Interrupt Control Register (\texttt{INTCTL}, Offset \texttt{0x0011})}
	\label{fig:reg-intctl}
\end{figure}

\begin{table}[H]
	\centering
	\begin{tabular}{llp{7.5cm}}
		\toprule
		\textbf{Bits} & \textbf{Field} & \textbf{Meaning}                                                                                           \\
		\midrule
		15--2         & ---            & Reserved                                                                                                   \\
		1             & \texttt{LIE}   & Line interrupt enable: when 1, the line-compare interrupt can be raised (Section~\ref{sec:line-interrupt}) \\
		0             & \texttt{VIE}   & Vertical sync interrupt enable: when 1, the device raises the NMI at the start of vertical blanking        \\
		\bottomrule
	\end{tabular}
	\caption{Interrupt Control Register Fields}
	\label{tab:reg-intctl-fields}
\end{table}

Both bits are 0 at reset, so no interrupt is raised until the program has set up its handlers.

\section{Reset State}
\label{sec:reset}

At reset the device is blank and every register is zero, except as follows:

\begin{table}[H]
	\centering
	\begin{tabular}{lll}
		\toprule
		\textbf{Register} & \textbf{Value at reset} & \textbf{Meaning}                                                      \\
		\midrule
		\texttt{CONFIG}   & \texttt{0x8078}         & \texttt{GD} = 1 (picture blank), \texttt{BRT} = 15, \texttt{MODE} = 0 \\
		\bottomrule
	\end{tabular}
	\caption{Registers with a Non-Zero Reset Value}
	\label{tab:reset}
\end{table}

The contents of video memory, the palette and the sprite table are undefined at reset. A program must initialise
whatever it uses, and then clear \texttt{GD}, before the picture is shown. A program should hide every sprite it does
not use (layer 3) before it clears \texttt{GD}.

